\chapter{Transmisión de  ISDB-Tb}


	En Argentina, la transmisión de TDT está regulada por el ENACOM a través de la resolución 7/13. Esta establece el esquema ISDB-Tb definido por la norma ABNT NBR 15601. Además, en el Anexo 1 se establece las especificaciones técnicas, tales como el tratamiento de datos, la modulación y la transmisión. En este documento también se definen los modos de recepción, especificándose la recepción total o parcial, denominada 1-seg. Asimismo, la normativa describe cómo se divide el espectro de un canal de 8 MHz en 13 segmentos, siendo el segmento central el destinado a la recepción parcial.


	La Figura \ref{fig:bloques_isdb} presenta el diagrama de bloques del transmisor de TDT. Este capítulo se encarga de explicar cada uno de ellos en el contexto de transmisión y recepción parcial 1-seg. Dado que en este trabajo se utiliza exclusivamente la capa jerárquica A, no se desarrollarán los bloques asociados a la división y combinación jerárquica.

	\figura{.9}{./isdb/orden_canales.png}{Orden de los segmentos en el canal de 8 MHz.}{ Fuente: \cite{isdbt}.}{\label{fig:canales_isdb}}

	\figura{.9}{./isdb/1seg.png}{Recepción total y parcial.}{ Fuente: \cite{isdbt}.}{\label{fig:1seg}}


	\figura{.9}{./isdb/cod_bloques.png}{Codificación de canal.}{ Fuente: \cite{isdbt}}{\label{fig:bloques_isdb}}





	\section{Código Reed-Solomon}
	Inicialmente, se aplicará la codificación externa en la que se emplea un código Reed-Solomon para agregar redundancia a los datos de entrada. En ISDB-T se emplea un código RS(204,188), que agrega 16 bytes de redundancia a cada bloque de 188 bytes de información utilizando los siguientes polinomios.
	\begin{equation}
	\text{Polinomio Generador: } g(x) = (x + \lambda^0)(x + \lambda^1)(x + \lambda^2) \cdots (x + \lambda^{15})
	\end{equation}

	\begin{equation}
		\text{Polinomio de Galois: }p(x) = x^8 + x^4 + x^3 + x^2 + 1
	\end{equation}
	donde $\lambda$ es el primo 2 del campo finito GF($2^8$) definido. Este esquema permite corregir hasta ocho bytes erróneos por bloque.


	\section{Dispersión de energía}

	Luego de la codificación inicial de los \textit{transport streams}, los datos ingresan al bloque de dispersión de energía. Esta etapa tiene como finalidad evitar la presencia de patrones repetitivos prolongados en la secuencia binaria, los cuales podrían concentrar energía en determinadas componentes espectrales.

	La dispersión se realiza utilizando una compuerta XOR para combinar los bits de entrada con una PSBR generada mediante un LSFR definido por un estado inicial \textit{100101010000000} y por el polinomio

	\begin{equation}
		G(x) = x^{15} + x^{14} + 1
	\end{equation}

	\section{Entrelazador de byte}
	Luego de que los datos sean protegidos por medio de la codificación RS y por la dispersión de energía, se entrelazan a nivel de byte mediante un interleaver convolucional. Este consiste en una cascada de colas FIFO con distintos niveles de profundidad, tal como se muestra en la Figura \ref{fig:intrlv_conv}, de manera que las distintas ramas sufran distintos retardos y así generar el entrelazaminento.


	\figura{.9}{./isdb/intrlv_conv.png}{Entrelazador de bytes.}{ Adaptado de \cite{isdbt}.}{\label{fig:intrlv_conv}}


	\section{Código convolucional}

	El segundo nivel de codificación, o codificación interna, se basa en un código convolucional con tasa base $1/2$,
	definido por los polinomios generadores

	\begin{equation}
		G_1 = 111_2, \qquad G_2 = 101_2 .
	\end{equation}

	A partir de este codificador se obtienen distintas tasas efectivas mediante el uso de
	patrones de puntuado. La elección del patrón depende del compromiso deseado entre
	robustez frente a errores y eficiencia espectral, y varía según el modo de transmisión
	configurado.



	\figura{.9}{./isdb/conv_code.png}{Diagrama en bloques del código convolucional.}{ Fuente: \cite{isdbt}.}{\label{fig:conv_code}}

	\section{Entrelazador de bit}
	Primeramente, se toma el bit-stream resultante de la codificación convolucional y se lo paraleliza en grupos de $log_2(M)$ bits según la modulación seleccionada siguiendo, siendo M el orden de la constelacion. Luego, se aplica un segundo nivel de entrelazado, esta vez a nivel de bit, utilizando un esquema similar al del entrelazador de bytes pero con registros que almacenan un solo bit. La configuración de las longitudes de los buffers en este caso depende del tipo de modulación seleccionada (QPSK, 16-QAM o 64-QAM) y se muestra en las Figuras \ref{fig:bit_intrl_4}, \ref{fig:bit_intrl_16} y \ref{fig:bit_intrl_64}.

	\figura{.9}{./isdb/bit_intrl_4.png}{Entralazador de bit QPSK.}{ Adaptado de \cite{isdbt}.}{\label{fig:bit_intrl_4}}

	\figura{.9}{./isdb/bit_intrl_16.png}{Entralazador de bit 16 QAM.}{ Adaptado de \cite{isdbt}.}{\label{fig:bit_intrl_16}}

	\figura{.9}{./isdb/bit_intrl_64.png}{Entralazador de bit 64QAM.}{ Adaptado de \cite{isdbt}.}{\label{fig:bit_intrl_64}}


	\section{Mapeo de símbolos QAM}

	El mapeo de los bits se realiza asignando los bits pares y los bits impares a las componentes real (I) e imaginaria (Q), respectivamente, siguiendo LMSB para el orden de los bits. A continuación, se detallan las constelaciones y mapeos correspondientes a cada tipo de modulación.
	\subsubsection{QPSK}
	\figura{.9}{./isdb/qpsk.png}{Mapeo de bits a símbolos QPSK.}{ Adaptado de \cite{isdbt}.}{\label{fig:qpsk}}
	\begin{equation}
		a_k = 1 - 2  b_0  +j (1 - 2 b_1)
	\end{equation}


	\subsubsection{16 QAM}
	\figura{.9}{./isdb/16qam.png}{Mapeo de bits a símbolos 16 QAM.}{ Adaptado de \cite{isdbt}.}{\label{fig:16qam}}

	\begin{align}
		\mathbb{R}\{a_k\} &= (1 - 2 b_0)\left[1 + 2(1 - 2 b_2)\right] \\
		\mathbb{I}\{a_k\} &= (1 - 2 b_1)\left[1 + 2(1 - 2 b_3)\right]
	\end{align}



	\subsubsection{64 QAM}
	\figura{.9}{./isdb/64qam.png}{Mapeo de bits a símbolos 64 QAM.}{ Adaptado de \cite{isdbt}.}{\label{fig:64qam}}
	\begin{align}
		\mathbb{R}\{a_k\} &= (1 - 2 b_0) \left[ 1 + 2(1 - 2 b_2) + 4(1 - 2 b_4) \right] \\
		\mathbb{I}\{a_k\} &= (1 - 2 b_1) \left[ 1 + 2(1 - 2 b_3) + 4(1 - 2 b_5) \right]
	\end{align}

	\subsection{Entrelazador de tiempo y frecuencia}

	Además del entrelazado a nivel de bit y de byte, el sistema usa dos tipos más de entrelazamiento, de tiempo y frecuencia. Estos bloques distribuyen los símbolos modulados tanto entre distintos símbolos OFDM (tiempo) como entre diferentes subportadoras (frecuencia).

	\subsubsection{Entrelazador de tiempo}

	Nuevamente, el entrelazador de tiempo está compuesto por una cascada de $n_c$ buffers FIFO de longitudes variables. Este esquema permite aplicarle distintos retardos a cada símbolo y así mezclándolos entre distintos símbolos OFDM.

	La longitud $L_i$ del $i$-ésimo buffer está dada por la siguiente expresión:

	\begin{equation}
		L_i = I \cdot m_i, \quad 0 \leq i \leq n_c - 1
	\end{equation}

	donde

	\begin{equation}
		m_i = i \cdot 5 \mod 96, \quad 0 \leq i \leq n_c - 1
	\end{equation}
	y siendo $n_c$ igual a 96, 192 y 384 en los modos 1, 2 3, respectivamente. La constante $I$ es un parámetro de transmisión configurable, que puede asumir uno de los cuatro valores indicados en la Tabla~\ref{tab:parametro_I} para un modo de transmisión determinado. La Figura~\ref{fig:intercalador_tiempo} muestra un diagrama en bloques en donde se puede observar la cascada de buffers y la configuración de sus longitudes.


	\figura{.9}{./isdb/time_intrl.png}{Entrelazador de tiempo}{ Fuente: \cite{isdbt}.}{\label{fig:intercalador_tiempo}}

	\begin{table}[h]
		\centering

		\renewcommand{\arraystretch}{1.2} % aumenta el espaciado vertical
		\begin{tabular}{|c c|}
			\hline
			\textbf{Modo} & \textbf{Valores de $I$} \\
			\hline\hline
			Modo 1 & 0, 4, 8, 16 \\
			Modo 2 & 0, 2, 4, 8 \\
			Modo 3 & 0, 1, 2, 3 \\
			\hline
		\end{tabular}
		\caption{Valores posibles del parámetro $I$ según el modo de transmisión. Adaptado de \cite{isdbt}.}
		\label{tab:parametro_I}
	\end{table}

	\subsubsection{Entrelazador de frecuencia}
	El entrelazador de frecuencia se encarga de distribuir los símbolos modulados entre las distintas subportadoras del frame OFDM. Para esto se utilizan dos esquemas, una primera instancia de rotación de portadoras y una segunda instancia de permutación de portadoras.

	Para la rotación de portadoras, dependiendo del modo en que se esté transmitiendo, se rotan las portadoras siguiendo los siguientes patrones
	\begin{table}[H]
		\centering
		\begin{tabular}{|c c|}
			\hline
			\textbf{Modo} & \textbf{Rotación de portadoras} \\
			\hline\hline
			Modo 1 & $i = (i + k) \mod 96$ \\
			Modo 2 & $i = (i + k) \mod 192$ \\
			Modo 3 & $i = (i + k) \mod 384$ \\
			\hline
		\end{tabular}
		\caption{Rotación de portadoras según el modo de transmisión. Adaptado de \cite{isdbt}.}
		\label{tab:rotacion_portadoras}
	\end{table}

	donde \textit{i} representa el índice de la portadora y \textit{k} representa la cantidad de segmentos de datos que se hayan transmitido hasta el momento. Luego de la rotación, se aplican las siguientes permutaciones de portadoras dependiendo del modo de transmisión:

	\begin{table}[H]
		\centering
		\begin{tabular}{|c c|}
			\hline
			Antes & Después \\\hline\hline
			&\\\hline
		\end{tabular}
		\caption{Permutación de portadoras para el Modo 1. Adaptado de \cite{isdbt}.}
		\label{tab:perm_m1}
	\end{table}

	\begin{table}[H]
		\centering
		\begin{tabular}{|c c|}
			\hline
			Antes & Después \\\hline\hline
			&\\\hline

		\end{tabular}
		\caption{Permutación de portadoras para el Modo 2. Adaptado de \cite{isdbt}.}
		\label{tab:perm_m2}
	\end{table}

	\begin{table}[H]
		\centering
		\begin{tabular}{|c c|}
			\hline
			Antes & Después \\\hline\hline

			&\\\hline
		\end{tabular}
		\caption{Permutación de portadoras para el Modo 3. Adaptado de \cite{isdbt}.}
		\label{tab:perm_m3}
	\end{table}



	\section{Estructura de frame OFDM}
	Una vez obtenidos los datos finales, es necesario estructurarlos dentro de un frame OFDM. Para esto se crean data frames que dependiendo del modo tendrán una longitud de 96, 192 y 384 portadoras. Estos data frame se agruparan en un arreglo bidimensional; como lo muestra la Figura \ref{fig:frame} en donde las filas representan el dominio del tiempo y las columnas el dominio de la frecuencia.


	\figura{.9}{./isdb/frame.png}{Estructura de datos en el frame OFDM.}{ Fuente: \cite{isdbt}.}{\label{fig:frame}}
	donde $a_{i,j,k}$ representa el késimo segmento del símbolo de la portadora, siendo $"$i$"$ la dirección de la portadora en el segmento OFDM y $"$j$"$ la dirección del símbolo en el segmento OFDM.


	\subsection{Secuencia pseudoaleatoria}

	Dentro del frame OFDM, además de las portadoras de datos, se incluyen otros tipos de portadoras auxiliares. En particular, algunas de ellas utilizan una secuencia pseudoaleatoria para la definición de sus símbolos.

	La secuencia pseudoaleatoria, denotada como $w$, se genera mediante un registro de desplazamiento con realimentación lineal (LFSR), definido por el polinomio generador

	\begin{equation}
		g(x) = 1 + x^{9} + x^{11}.
	\end{equation}

	Para el segmento central del frame, la secuencia se inicializa utilizando la seed indicada en la Tabla []. En la Figura~\ref{fig:prbs} se muestra el esquema del lfsr para la generación de dicha secuencia.


	\figura{.9}{./isdb/prbs.png}{Generador de secuencia binaria pseudo aleatoria .}{ Fuente: \cite{isdbt}.}{\label{fig:prbs}}


	\subsection{Señales SP (scattered pilots)}

	Las portadoras SP se encuentran distribuidas de forma dispersa dentro del frame, presentando periodicidad tanto dentro de un mismo símbolo OFDM como entre símbolos OFDM consecutivos, de acuerdo con el patrón

	 \begin{equation}
	 		\left\{ i = I_{min} + 3 \cdot (j mod 4) + 12 p \hspace{5pt}|\hspace{5pt} p \in \mathbb{N},\hspace{5pt} i \in [I_{min}; I_{max}] \right\}
	 \end{equation}

	 Quedando estas periodicas cada 12 portadoras dentro del simbolo y cada 4 portadoras entre simbolos. Estas portadoras se utilizan para la estimación de la respuesta en frecuencia del canal $H(t,f)$ en un determinado instante de tiempo, lo que permite compensar los efectos introducidos por el canal sobre la señal transmitida.

	Las portadoras SP se modulan mediante BPSK. No obstante, el valor transmitido en cada una de ellas no corresponde a un dato de información, sino al valor de una secuencia pseudoaleatoria asociada a la frecuencia específica de la portadora.

	Una vez obtenida la secuencia $w$, en las portadoras indicadas en la Tabla~[X] se asignan los símbolos

	{\Huge[Completar tablar portadoras SP]}

	\begin{equation}
		a_{i,j,k} = \frac{4}{3} - \frac{8}{3} w_i,
	\end{equation}

	donde $i$ representa el índice de la portadora. De esta forma, los valores binarios $w_i \in \{0,1\}$ se mapean en los niveles $\left\{ \frac{4}{3}, -\frac{4}{3} \right\}$, respectivamente.

	\figura{.9}{./isdb/pilotos.png}{Distribucion de portadoras pilotos.}{ Fuente: \cite{isdbt}.}{\label{fig:pilotos}}

	\subsection{Señales AC (auxiliary channel)}
	Las portadoras AC se utilizan para la transmisión de información auxiliar y de control necesaria para la programación del sistema. Esta información se modula mediante DBPSK a lo largo de 203 de los 204 símbolos que componen el superframe OFDM.

	{\Huge[Completar tablar portadoras AC]}


	En el primer símbolo del superframe, las portadoras indicadas en la Tabla~[X] se modulan utilizando BPSK, de la misma forma que las portadoras SP. A partir del segundo símbolo, se consideran los 203 bits correspondientes a la información AC y se aplica un esquema de codificación diferencial entre símbolos consecutivos.

	La operación diferencial se define como

	\begin{equation}
		a_{i,j,k} = a_{i-1,j,k} \oplus b_i
	\end{equation}

	De este modo, las portadoras AC del símbolo actual se invierten o se mantienen iguales respecto al símbolo anterior para representar un valor lógico 1 o 0, respectivamente.


	\subsection{Señal TMCC}
	{\Huge[Completar tablar portadoras TMCC]}
	\section{Frecuencia de muestreo y frecuencias centrales}

% -------------------------------------------------------------------------
% Capítulo 4 — Recepción ISDB-Tb (cadena teórica)
% Basado en la cadena de transmisión del Capítulo 3 del documento provisto.
% :contentReference[oaicite:0]{index=0}
% -------------------------------------------------------------------------

\chapter{Recepción}
\label{ch:recepcion_isdbt}

En este capítulo se describe la cadena de recepción de una señal ISDB-Tb,
tomando como referencia la estructura de transmisión presentada en el Capítulo~3
y describiendo las etapas necesarias para recuperar los datos y parámetros de transmisión.
La recepción se presenta como una secuencia de bloques funcionales, desde la señal
compleja en banda base hasta la reconstrucción del \textit{Transport Stream} (TS).

\section{Diagrama general de la cadena de recepción}
\label{sec:rx_diagrama_general}

A nivel de alto nivel, la cadena de recepción puede resumirse como:

\begin{center}
	\begin{minipage}{0.95\textwidth}
		\begin{verbatim}
			IQ (banda base)
			-> Sincronización temporal (detección de símbolo) y corrección de CFO
			-> Remoción de prefijo cíclico + FFT (demodulación OFDM)
			-> Sincronización fina / alineamiento de portadoras con pilotos (SP)
			-> Estimación de canal y ecualización
			-> Extracción y decodificación de TMCC/AC (parámetros de transmisión)
			-> Demapeo QAM (LLR/bits)
			-> De-entrelazado (tiempo/frecuencia, bit, byte)
			-> Depuncturing + Decodificación Viterbi (código convolucional)
			-> Decodificación Reed-Solomon
			-> Des-dispersión de energía (descrambling)
			-> TS recuperado
		\end{verbatim}
	\end{minipage}
\end{center}

\section{Etapas de sincronización y demodulación OFDM}
\label{sec:rx_ofdm_sync}

\subsection{Adquisición y conversión a banda base}
La señal recibida se considera como una secuencia compleja IQ muestreada a la frecuencia
correspondiente al canal/segmento de interés. En el caso de recepción parcial (1-seg),
la señal se limita al segmento central y se procesa en banda base compleja.

\subsection{Sincronización temporal: detección del inicio de símbolo}
La primera etapa crítica es la identificación del comienzo de los símbolos OFDM.
Aprovechando el prefijo cíclico (CP), se puede obtener una métrica de correlación entre
el CP y la porción correspondiente del final del símbolo. El máximo de esta métrica
permite estimar el instante de inicio del símbolo, y por ende segmentar la señal en
bloques símbolo a símbolo.

\subsection{Corrección de error de frecuencia (CFO) y fase}
La diferencia de fase entre el CP y su réplica en el final del símbolo permite estimar
un corrimiento de frecuencia (CFO) y/o una rotación de fase. Con esa estimación se aplica
una corrección compleja de la forma $x[n]\exp(-j2\pi \widehat{\Delta f} n / f_s)$
para centrar el espectro antes de la FFT.

\subsection{Remoción de CP y FFT}
Una vez detectado el símbolo y compensado el CFO (al menos en forma gruesa),
se elimina el prefijo cíclico y se calcula la FFT para obtener las subportadoras en el
dominio de la frecuencia. Esto produce, por símbolo, el vector de portadoras que contiene:
portadoras de datos, pilotos dispersos (SP) y portadoras de señalización (AC/TMCC).

\section{Pilotos, estimación de canal y ecualización}
\label{sec:rx_channel_est}

\subsection{Alineamiento fino con pilotos dispersos (SP)}
Las portadoras SP ocupan posiciones predeterminadas y transmiten una secuencia conocida.
Dicha propiedad permite realizar sincronización fina (por ejemplo, corrección de
desplazamientos residuales y verificación de alineamiento de índices de portadora)
mediante correlación o búsqueda de la mejor coincidencia entre el patrón esperado y el
observado.

\subsection{Estimación de canal}
Con los símbolos pilotos (SP) se estima la respuesta compleja del canal
$H[k]$ (por símbolo OFDM y por portadora). Luego se interpola la estimación
en frecuencia y/o tiempo para obtener una estimación del canal en las portadoras de datos.

\subsection{Ecualización}
Con $\widehat{H}[k]$, se ecualiza cada portadora de datos:
\[
\widehat{X}[k] = \frac{Y[k]}{\widehat{H}[k]},
\]
donde $Y[k]$ es la portadora recibida y $\widehat{X}[k]$ la portadora ecualizada
lista para demapeo.

\section{Extracción de información de control: TMCC y AC}
\label{sec:rx_tmcc_ac}

Además de los datos, el receptor debe recuperar los parámetros de transmisión
(modulación, tasas de codificación, entrelazados, etc.). Para ello:
\begin{itemize}
	\item Se extraen las portadoras correspondientes a AC y TMCC, ubicadas en índices
	definidos por la norma.
	\item Se demodulan según su esquema (típicamente DBPSK para señalización),
	y se decodifica la información de configuración.
\end{itemize}
A partir de la decodificación de TMCC/AC, el receptor conoce (o verifica) la configuración
necesaria para los bloques posteriores: tipo de constelación, tasa del código convolucional
(patrón de puntuado), y parámetros de entrelazado temporal.

\section{Demapeo y cadena de decodificación de canal (inversa a transmisión)}
\label{sec:rx_fec_chain}

A continuación se describen los bloques que invierten conceptualmente la codificación
realizada en transmisión.

\subsection{Demapeo de símbolos QAM}
Las portadoras ecualizadas se demapean a bits (o métricas tipo LLR) de acuerdo con la
constelación configurada (QPSK/16QAM/64QAM). Esta etapa produce una secuencia binaria
que alimenta la decodificación FEC.

\subsection{De-entrelazado en tiempo/frecuencia}
Dado que en transmisión se aplica un entrelazado adicional en los dominios temporal y
frecuencial, el receptor aplica el proceso inverso para reordenar los símbolos hacia su
posición original, reduciendo así el efecto de errores en ráfaga y desvanecimientos selectivos.

\subsection{De-entrelazado de bit}
Se aplica el de-entrelazado a nivel de bit (inverso del entrelazador convolucional de bit)
según la modulación empleada. El objetivo es restaurar el orden original de los bits antes
del FEC interno.

\subsection{Depuncturing y decodificación convolucional (Viterbi)}
En transmisión, el código convolucional base se combina con un patrón de puntuado para
obtener distintas tasas efectivas. En recepción:
\begin{itemize}
	\item Se aplica \textit{depuncturing} para reconstruir la secuencia equivalente a tasa $1/2$,
	insertando marcadores de borrado donde corresponda.
	\item Se ejecuta el decodificador Viterbi para recuperar la secuencia de bits más probable.
\end{itemize}

\subsection{De-entrelazado de byte}
Se revierte el entrelazado de bytes aplicado en transmisión, reordenando los bytes para
que los errores vuelvan a concentrarse y puedan ser corregidos por la etapa RS.

\subsection{Decodificación Reed-Solomon}
Se aplica el decodificador Reed-Solomon RS(204,188), eliminando los bytes de paridad e
intentando corregir errores (particularmente en ráfaga) dentro de cada bloque.

\subsection{Des-dispersión de energía (descrambling)}
Finalmente, se invierte la dispersión de energía mediante la misma secuencia pseudoaleatoria
(LFSR) utilizada en transmisión, aplicando nuevamente XOR para recuperar la secuencia original.

\section{Salida}
\label{sec:rx_output}

El resultado del proceso es el \textit{Transport Stream} (TS) recuperado, listo para ser utilizado
en etapas posteriores del sistema (por ejemplo, reconstrucción de la señal transmitida para
canal de referencia coherente en radar pasivo, o para demultiplexado y análisis de contenido).
ss
% -------------------------------------------------------------------------
% Nota práctica de redacción:
% Si querés mantener simetría con el Capítulo 3, podés usar las mismas
% subsecciones (3.1..3.7) pero con prefijo "De-" en recepción.
% -------------------------------------------------------------------------
